% C3E Revision R1 methods supplement. All tables are generated by
% `python -m c3e.revision.report` from results/c3e_revision/.
\documentclass[10pt,a4paper]{article}
\usepackage[utf8]{inputenc}
\usepackage[T1]{fontenc}
\usepackage{lmodern}
\usepackage[margin=0.8in]{geometry}
\usepackage{amsmath,amssymb}
\usepackage{booktabs,longtable,array,adjustbox}
\usepackage{microtype}
\usepackage{pdflscape}
\usepackage[hidelinks]{hyperref}
\usepackage[capitalise,noabbrev]{cleveref}
\setlength{\emergencystretch}{3em}
% Generated by c3e.revision.report from results/c3e_revision; do not edit.
\newcommand{\APShareExt}{83\%}
\newcommand{\APShareSrc}{60\%}
\newcommand{\AgeCovDispMFour}{0.033}
\newcommand{\AgeCovDispMThree}{0.115~(+0.105, +0.125)}
\newcommand{\AgreeMaxExt}{99.7\%}
\newcommand{\AgreeMaxSrc}{98.6\%}
\newcommand{\AgreeMinExt}{97.7\%}
\newcommand{\AgreeMinSrc}{95.2\%}
\newcommand{\CTwoIdentExt}{34}
\newcommand{\CTwoIdentSrc}{6}
\newcommand{\CTwoMultiDonorExt}{2,564}
\newcommand{\CTwoMultiDonorSrc}{1,465}
\newcommand{\CTwoMultiImgExt}{2,608}
\newcommand{\CTwoMultiImgSrc}{1,607}
\newcommand{\CTwoRowsExt}{135,561}
\newcommand{\CTwoRowsSrc}{16,456}
\newcommand{\CTwoSamePatExt}{0}
\newcommand{\CTwoSamePatSrc}{0}
\newcommand{\CTwoTermDiffExt}{28\%}
\newcommand{\CTwoTermDiffSrc}{54\%}
\newcommand{\CTwoTokensExt}{7.3}
\newcommand{\CTwoTokensSrc}{13.0}
\newcommand{\CalRatioMax}{3.4}
\newcommand{\CalRatioMin}{1.5}
\newcommand{\CalSDMFour}{0.010}
\newcommand{\ConfAccMax}{59\%}
\newcommand{\ConfAccMin}{47\%}
\newcommand{\CovDriftMaxCZero}{0.047}
\newcommand{\DensExtFindings}{9.1\%}
\newcommand{\DensExtImpression}{28.8\%}
\newcommand{\DensSrcFindings}{44.6\%}
\newcommand{\DensSrcImpression}{20.6\%}
\newcommand{\ELMFourExtCov}{0.855}
\newcommand{\EvalSDMFour}{0.004}
\newcommand{\ExtAtelPos}{19,431}
\newcommand{\ExtAtelPosRate}{97.9\%}
\newcommand{\ExtAtelUnc}{21,462}
\newcommand{\ExtEdemaUnmNeg}{25.5\%}
\newcommand{\ExtEffUnmNeg}{43.5\%}
\newcommand{\ExtImages}{191,070}
\newcommand{\ExtNOne}{132,923}
\newcommand{\ExtNOneImages}{135,561}
\newcommand{\ExtNOnePatients}{53,110}
\newcommand{\ExtNThree}{40,622}
\newcommand{\ExtNTwo}{14,129}
\newcommand{\ExtPatients}{64,709}
\newcommand{\ExtStudies}{187,673}
\newcommand{\FailAUROCELMax}{0.76}
\newcommand{\FailAUROCELMin}{0.69}
\newcommand{\FailAUROCFrozenMax}{0.65}
\newcommand{\FailAUROCFrozenMin}{0.60}
\newcommand{\HFourMFourexternalEval}{$-$0.0077~($-$0.0101, $-$0.0054)}
\newcommand{\HFourMFourexternalEvalPt}{$-$0.0077}
\newcommand{\HFourMFourexternalOrig}{$-$0.0060~($-$0.0080, $-$0.0041)}
\newcommand{\HFourMFourexternalOrigPt}{$-$0.0060}
\newcommand{\HFourMFoursourceEval}{+0.0301~(+0.0215, +0.0383)}
\newcommand{\HFourMFoursourceEvalPt}{+0.0301}
\newcommand{\HFourMFoursourceOrig}{+0.0173~(+0.0125, +0.0217)}
\newcommand{\HFourMFoursourceOrigPt}{+0.0173}
\newcommand{\HFourMThreeexternalEval}{+0.0263~(+0.0240, +0.0286)}
\newcommand{\HFourMThreeexternalEvalPt}{+0.0263}
\newcommand{\HFourMThreeexternalOrig}{+0.0231~(+0.0211, +0.0250)}
\newcommand{\HFourMThreeexternalOrigPt}{+0.0231}
\newcommand{\HFourMThreesourceEval}{+0.0783~(+0.0691, +0.0865)}
\newcommand{\HFourMThreesourceEvalPt}{+0.0783}
\newcommand{\HFourMThreesourceOrig}{+0.0475~(+0.0422, +0.0524)}
\newcommand{\HFourMThreesourceOrigPt}{+0.0475}
\newcommand{\HOneMFourEval}{$-$0.0018~($-$0.0108, +0.0064)}
\newcommand{\HOneMFourEvalPt}{$-$0.0018}
\newcommand{\HOneMFourOrig}{+0.0384~(+0.0336, +0.0430)}
\newcommand{\HOneMFourOrigPt}{+0.0384}
\newcommand{\HOneMOneEval}{+0.0120~($-$0.0002, +0.0244)}
\newcommand{\HOneMOneEvalPt}{+0.0120}
\newcommand{\HOneMOneOrig}{+0.0734~(+0.0668, +0.0802)}
\newcommand{\HOneMOneOrigPt}{+0.0734}
\newcommand{\HOneMThreeEval}{$-$0.0042~($-$0.0145, +0.0061)}
\newcommand{\HOneMThreeEvalPt}{$-$0.0042}
\newcommand{\HOneMThreeOrig}{+0.0469~(+0.0408, +0.0531)}
\newcommand{\HOneMThreeOrigPt}{+0.0469}
\newcommand{\HOneMTwoEval}{$-$0.0173~($-$0.0292, $-$0.0054)}
\newcommand{\HOneMTwoEvalPt}{$-$0.0173}
\newcommand{\HOneMTwoOrig}{+0.0763~(+0.0687, +0.0838)}
\newcommand{\HOneMTwoOrigPt}{+0.0763}
\newcommand{\HThreeMFourEval}{$-$0.0138~($-$0.0243, $-$0.0029)}
\newcommand{\HThreeMFourEvalPt}{$-$0.0138}
\newcommand{\HThreeMFourOrig}{$-$0.0350~($-$0.0408, $-$0.0292)}
\newcommand{\HThreeMFourOrigPt}{$-$0.0350}
\newcommand{\HThreeMThreeEval}{$-$0.0162~($-$0.0256, $-$0.0070)}
\newcommand{\HThreeMThreeEvalPt}{$-$0.0162}
\newcommand{\HThreeMThreeOrig}{$-$0.0265~($-$0.0321, $-$0.0212)}
\newcommand{\HThreeMThreeOrigPt}{$-$0.0265}
\newcommand{\HTwoMFourCOneEval}{+0.0089~(+0.0002, +0.0173)}
\newcommand{\HTwoMFourCOneEvalPt}{+0.0089}
\newcommand{\HTwoMFourCOneOrig}{+0.0096~(+0.0048, +0.0144)}
\newcommand{\HTwoMFourCOneOrigPt}{+0.0096}
\newcommand{\HTwoMFourCTwoEval}{$-$0.0289~($-$0.0377, $-$0.0199)}
\newcommand{\HTwoMFourCTwoEvalPt}{$-$0.0289}
\newcommand{\HTwoMFourCTwoOrig}{$-$0.0137~($-$0.0185, $-$0.0087)}
\newcommand{\HTwoMFourCTwoOrigPt}{$-$0.0137}
\newcommand{\HTwoMThreeCOneEval}{+0.0048~($-$0.0038, +0.0134)}
\newcommand{\HTwoMThreeCOneEvalPt}{+0.0048}
\newcommand{\HTwoMThreeCOneOrig}{+0.0024~($-$0.0027, +0.0076)}
\newcommand{\HTwoMThreeCOneOrigPt}{+0.0024}
\newcommand{\HTwoMThreeCTwoEval}{$-$0.0472~($-$0.0572, $-$0.0374)}
\newcommand{\HTwoMThreeCTwoEvalPt}{$-$0.0472}
\newcommand{\HTwoMThreeCTwoOrig}{$-$0.0219~($-$0.0279, $-$0.0160)}
\newcommand{\HTwoMThreeCTwoOrigPt}{$-$0.0219}
\newcommand{\InsCovDispMFour}{0.019}
\newcommand{\JaccardMax}{0.85}
\newcommand{\JaccardMin}{0.64}
\newcommand{\LSAOneMFourEval}{$-$0.0389~($-$0.0480, $-$0.0297)}
\newcommand{\LSAOneMFourOrig}{$-$0.0839~($-$0.0899, $-$0.0781)}
\newcommand{\LSAOneMThreeEval}{+0.0090~($-$0.0001, +0.0185)}
\newcommand{\LSAOneMThreeOrig}{$-$0.0613~($-$0.0667, $-$0.0557)}
\newcommand{\LSAOnecFMFourEval}{$-$0.0271~($-$0.0400, $-$0.0142)}
\newcommand{\LSAOnecFMFourOrig}{$-$0.0077~($-$0.0104, $-$0.0050)}
\newcommand{\LSAOnecFMThreeEval}{$-$0.0090~($-$0.0257, +0.0094)}
\newcommand{\LSAOnecFMThreeOrig}{$-$0.0069~($-$0.0107, $-$0.0030)}
\newcommand{\LSAOnecIMFourEval}{$-$0.0212~($-$0.0347, $-$0.0088)}
\newcommand{\LSAOnecIMFourOrig}{$-$0.0065~($-$0.0091, $-$0.0040)}
\newcommand{\LSAOnecIMThreeEval}{$-$0.0042~($-$0.0212, +0.0126)}
\newcommand{\LSAOnecIMThreeOrig}{$-$0.0059~($-$0.0096, $-$0.0022)}
\newcommand{\LSAThreeMFourEval}{+0.0048~($-$0.0034, +0.0128)}
\newcommand{\LSAThreeMFourOrig}{$-$0.1079~($-$0.1141, $-$0.1014)}
\newcommand{\LSAThreeMThreeEval}{+0.0285~(+0.0192, +0.0382)}
\newcommand{\LSAThreeMThreeOrig}{$-$0.1324~($-$0.1398, $-$0.1246)}
\newcommand{\LSATwoMFourEval}{$-$0.0389~($-$0.0478, $-$0.0299)}
\newcommand{\LSATwoMFourOrig}{$-$0.1610~($-$0.1682, $-$0.1535)}
\newcommand{\LSATwoMThreeEval}{+0.0090~($-$0.0002, +0.0188)}
\newcommand{\LSATwoMThreeOrig}{$-$0.1359~($-$0.1435, $-$0.1280)}
\newcommand{\LSATwoxMFourEval}{+0.0126~($-$0.0004, +0.0244)}
\newcommand{\LSATwoxMFourOrig}{+0.0664~(+0.0614, +0.0714)}
\newcommand{\LSATwoxMThreeEval}{+0.0012~($-$0.0141, +0.0159)}
\newcommand{\LSATwoxMThreeOrig}{+0.0809~(+0.0741, +0.0876)}
\newcommand{\LSAZeroMFourEval}{$-$0.0018~($-$0.0108, +0.0064)}
\newcommand{\LSAZeroMFourOrig}{+0.0384~(+0.0336, +0.0430)}
\newcommand{\LSAZeroMThreeEval}{$-$0.0042~($-$0.0145, +0.0061)}
\newcommand{\LSAZeroMThreeOrig}{+0.0469~(+0.0408, +0.0531)}
\newcommand{\LSBTwoMFourEval}{+0.0048~($-$0.0033, +0.0132)}
\newcommand{\LSBTwoMFourOrig}{$-$0.0524~($-$0.0571, $-$0.0477)}
\newcommand{\LSBTwoMThreeEval}{+0.0285~(+0.0193, +0.0380)}
\newcommand{\LSBTwoMThreeOrig}{$-$0.0579~($-$0.0633, $-$0.0525)}
\newcommand{\MFourCOneCovExt}{0.764}
\newcommand{\MFourCardPosExt}{76\%}
\newcommand{\MFourCardPosSrc}{56\%}
\newcommand{\MFourCardSensExt}{0.97}
\newcommand{\MFourCardSensSrc}{0.89}
\newcommand{\MFourCardSpecExt}{0.44}
\newcommand{\MFourCardSpecSrc}{0.82}
\newcommand{\MFourEdemaSpecExt}{0.65}
\newcommand{\MFourEdemaSpecSrc}{0.84}
\newcommand{\MFourEffSpecExt}{0.64}
\newcommand{\MFourEffSpecSrc}{0.89}
\newcommand{\MThreeCTwoCovExt}{0.746}
\newcommand{\MThreeCTwoCovSrc}{0.728}
\newcommand{\MTwoCOneCovExt}{1.000}
\newcommand{\MTwoCOneCovSrc}{1.000}
\newcommand{\MTwoMThreeReplicateStatus}{completed}
\newcommand{\NoLabExtFindings}{75.7\%}
\newcommand{\NoLabExtImpression}{18.7\%}
\newcommand{\NoLabNExtFindings}{100,617}
\newcommand{\NoLabNExtImpression}{24,821}
\newcommand{\NoLabNSrcFindings}{553}
\newcommand{\NoLabNSrcImpression}{6,968}
\newcommand{\NoLabSrcFindings}{6.0\%}
\newcommand{\NoLabSrcImpression}{47.2\%}
\newcommand{\OrderGapMM}{0.026}
\newcommand{\OrderGapMTwo}{0.049}
\newcommand{\PathMFourAtel}{$-$0.026}
\newcommand{\PathMFourCard}{+0.053}
\newcommand{\PathMFourUnmMax}{+0.18}
\newcommand{\PathMFourUnmMin}{+0.03}
\newcommand{\PermHFourMThreeExtRange}{+0.0252 to +0.0263}
\newcommand{\PermHFourMThreeSrcRange}{+0.0779 to +0.0801}
\newcommand{\PosExtFindings}{68.5\%}
\newcommand{\PosExtImpression}{70.0\%}
\newcommand{\PosSrcFindings}{33.3\%}
\newcommand{\PosSrcImpression}{71.1\%}
\newcommand{\RepTwoHFourMThreeexternaleval}{+0.0011~($-$0.0010, +0.0031)}
\newcommand{\RepTwoHFourMThreeexternalorig}{+0.0023~(+0.0006, +0.0040)}
\newcommand{\RepTwoHFourMThreesourceeval}{+0.0565~(+0.0484, +0.0647)}
\newcommand{\RepTwoHFourMThreesourceorig}{+0.0356~(+0.0311, +0.0404)}
\newcommand{\RepTwoHOneMFoureval}{$-$0.0469~($-$0.0586, $-$0.0350)}
\newcommand{\RepTwoHOneMFourorig}{+0.0246~(+0.0182, +0.0315)}
\newcommand{\RepTwoHOneMOneeval}{$-$0.0082~($-$0.0194, +0.0027)}
\newcommand{\RepTwoHOneMOneorig}{+0.0511~(+0.0450, +0.0568)}
\newcommand{\RepTwoHOneMThreeeval}{$-$0.0402~($-$0.0503, $-$0.0298)}
\newcommand{\RepTwoHOneMThreeorig}{+0.0145~(+0.0086, +0.0208)}
\newcommand{\RepTwoHOneMTwoeval}{$-$0.0279~($-$0.0411, $-$0.0147)}
\newcommand{\RepTwoHOneMTwoorig}{+0.0669~(+0.0586, +0.0749)}
\newcommand{\RepTwoHThreeMThreeeval}{$-$0.0320~($-$0.0405, $-$0.0237)}
\newcommand{\RepTwoHThreeMThreeorig}{$-$0.0366~($-$0.0414, $-$0.0318)}
\newcommand{\RepTwoHTwoMThreeeval}{+0.0150~(+0.0072, +0.0228)}
\newcommand{\RepTwoHTwoMThreeorig}{+0.0154~(+0.0108, +0.0198)}
\newcommand{\RiskEvalExtMFour}{0.1374}
\newcommand{\RiskEvalExtMOne}{0.2230}
\newcommand{\RiskEvalExtMThree}{0.1895}
\newcommand{\RiskEvalExtMTwo}{0.2869}
\newcommand{\RiskEvalSrcMFour}{0.1391}
\newcommand{\RiskEvalSrcMOne}{0.2111}
\newcommand{\RiskEvalSrcMThree}{0.1937}
\newcommand{\RiskEvalSrcMTwo}{0.3042}
\newcommand{\RiskOrigExtMFour}{0.1128}
\newcommand{\RiskOrigExtMOne}{0.1818}
\newcommand{\RiskOrigExtMThree}{0.1571}
\newcommand{\RiskOrigExtMTwo}{0.2338}
\newcommand{\RiskOrigSrcMFour}{0.0744}
\newcommand{\RiskOrigSrcMOne}{0.1084}
\newcommand{\RiskOrigSrcMThree}{0.1102}
\newcommand{\RiskOrigSrcMTwo}{0.1575}
\newcommand{\SOneCellsChanged}{44,180--246,554}
\newcommand{\SOneMOneCells}{44,180}
\newcommand{\SOneMOneStudies}{1,633}
\newcommand{\SOneMaxProbDiff}{$1.2\times10^{-7}$}
\newcommand{\SOneStudiesChanged}{178--4,518}
\newcommand{\ScanComparExt}{0.06\%}
\newcommand{\ScanComparSrc}{1.7\%}
\newcommand{\ScanConclExt}{0.02\%}
\newcommand{\ScanConclSrc}{0.45\%}
\newcommand{\ScanHeaderExt}{0.00\%}
\newcommand{\ScanHeaderSrc}{0.00\%}
\newcommand{\ScanQueryExt}{16.0\%}
\newcommand{\ScanQuerySrc}{59.8\%}
\newcommand{\ScanTargetExt}{15.8\%}
\newcommand{\ScanTargetSrc}{35.6\%}
\newcommand{\SeedCovDriftMax}{0.037}
\newcommand{\SensCellHOneMFour}{+0.0076~($-$0.0004, +0.0154)}
\newcommand{\SensCellHOneMFourPt}{+0.0076}
\newcommand{\SensCellHOneMOne}{+0.0156~(+0.0048, +0.0269)}
\newcommand{\SensCellHOneMOnePt}{+0.0156}
\newcommand{\SensCellHOneMThree}{$-$0.0026~($-$0.0120, +0.0069)}
\newcommand{\SensCellHOneMThreePt}{$-$0.0026}
\newcommand{\SensCellHOneMTwo}{$-$0.0013~($-$0.0136, +0.0109)}
\newcommand{\SensCellHOneMTwoPt}{$-$0.0013}
\newcommand{\SensCellHThreeMFour}{$-$0.0080~($-$0.0167, +0.0009)}
\newcommand{\SensCellHThreeMFourPt}{$-$0.0080}
\newcommand{\SensCellHThreeMThree}{$-$0.0183~($-$0.0264, $-$0.0099)}
\newcommand{\SensCellHThreeMThreePt}{$-$0.0183}
\newcommand{\SensCovNinetyHOneMFour}{$-$0.0058~($-$0.0147, +0.0022)}
\newcommand{\SensCovNinetyHOneMFourPt}{$-$0.0058}
\newcommand{\SensCovNinetyHOneMOne}{+0.0170~(+0.0053, +0.0291)}
\newcommand{\SensCovNinetyHOneMOnePt}{+0.0170}
\newcommand{\SensCovNinetyHOneMThree}{$-$0.0010~($-$0.0109, +0.0093)}
\newcommand{\SensCovNinetyHOneMThreePt}{$-$0.0010}
\newcommand{\SensCovNinetyHOneMTwo}{$-$0.0180~($-$0.0297, $-$0.0067)}
\newcommand{\SensCovNinetyHOneMTwoPt}{$-$0.0180}
\newcommand{\SensCovNinetyHThreeMFour}{$-$0.0229~($-$0.0331, $-$0.0123)}
\newcommand{\SensCovNinetyHThreeMFourPt}{$-$0.0229}
\newcommand{\SensCovNinetyHThreeMThree}{$-$0.0181~($-$0.0266, $-$0.0094)}
\newcommand{\SensCovNinetyHThreeMThreePt}{$-$0.0181}
\newcommand{\SensCovSeventyHOneMFour}{+0.0029~($-$0.0058, +0.0110)}
\newcommand{\SensCovSeventyHOneMFourPt}{+0.0029}
\newcommand{\SensCovSeventyHOneMOne}{+0.0081~($-$0.0049, +0.0211)}
\newcommand{\SensCovSeventyHOneMOnePt}{+0.0081}
\newcommand{\SensCovSeventyHOneMThree}{$-$0.0069~($-$0.0170, +0.0034)}
\newcommand{\SensCovSeventyHOneMThreePt}{$-$0.0069}
\newcommand{\SensCovSeventyHOneMTwo}{$-$0.0123~($-$0.0249, +0.0001)}
\newcommand{\SensCovSeventyHOneMTwoPt}{$-$0.0123}
\newcommand{\SensCovSeventyHThreeMFour}{$-$0.0052~($-$0.0161, +0.0061)}
\newcommand{\SensCovSeventyHThreeMFourPt}{$-$0.0052}
\newcommand{\SensCovSeventyHThreeMThree}{$-$0.0150~($-$0.0253, $-$0.0048)}
\newcommand{\SensCovSeventyHThreeMThreePt}{$-$0.0150}
\newcommand{\SensEvalHOneMFour}{$-$0.0018~($-$0.0108, +0.0064)}
\newcommand{\SensEvalHOneMFourPt}{$-$0.0018}
\newcommand{\SensEvalHOneMOne}{+0.0120~($-$0.0002, +0.0244)}
\newcommand{\SensEvalHOneMOnePt}{+0.0120}
\newcommand{\SensEvalHOneMThree}{$-$0.0042~($-$0.0145, +0.0061)}
\newcommand{\SensEvalHOneMThreePt}{$-$0.0042}
\newcommand{\SensEvalHOneMTwo}{$-$0.0173~($-$0.0292, $-$0.0054)}
\newcommand{\SensEvalHOneMTwoPt}{$-$0.0173}
\newcommand{\SensEvalHThreeMFour}{$-$0.0138~($-$0.0243, $-$0.0029)}
\newcommand{\SensEvalHThreeMFourPt}{$-$0.0138}
\newcommand{\SensEvalHThreeMThree}{$-$0.0162~($-$0.0256, $-$0.0070)}
\newcommand{\SensEvalHThreeMThreePt}{$-$0.0162}
\newcommand{\SensFindHOneMFour}{+0.0048~($-$0.0034, +0.0128)}
\newcommand{\SensFindHOneMFourPt}{+0.0048}
\newcommand{\SensFindHOneMOne}{+0.0951~(+0.0865, +0.1031)}
\newcommand{\SensFindHOneMOnePt}{+0.0951}
\newcommand{\SensFindHOneMThree}{+0.0285~(+0.0192, +0.0382)}
\newcommand{\SensFindHOneMThreePt}{+0.0285}
\newcommand{\SensFindHOneMTwo}{$-$0.0002~($-$0.0107, +0.0103)}
\newcommand{\SensFindHOneMTwoPt}{$-$0.0002}
\newcommand{\SensFindHThreeMFour}{$-$0.0903~($-$0.0983, $-$0.0817)}
\newcommand{\SensFindHThreeMFourPt}{$-$0.0903}
\newcommand{\SensFindHThreeMThree}{$-$0.0666~($-$0.0743, $-$0.0581)}
\newcommand{\SensFindHThreeMThreePt}{$-$0.0666}
\newcommand{\SensFindOrigHOneMFour}{$-$0.1079~($-$0.1141, $-$0.1014)}
\newcommand{\SensFindOrigHOneMFourPt}{$-$0.1079}
\newcommand{\SensFindOrigHOneMOne}{$-$0.0672~($-$0.0736, $-$0.0607)}
\newcommand{\SensFindOrigHOneMOnePt}{$-$0.0672}
\newcommand{\SensFindOrigHOneMThree}{$-$0.1324~($-$0.1398, $-$0.1246)}
\newcommand{\SensFindOrigHOneMThreePt}{$-$0.1324}
\newcommand{\SensFindOrigHOneMTwo}{$-$0.2385~($-$0.2470, $-$0.2303)}
\newcommand{\SensFindOrigHOneMTwoPt}{$-$0.2385}
\newcommand{\SensFindOrigHThreeMFour}{$-$0.0408~($-$0.0473, $-$0.0341)}
\newcommand{\SensFindOrigHThreeMFourPt}{$-$0.0408}
\newcommand{\SensFindOrigHThreeMThree}{$-$0.0652~($-$0.0713, $-$0.0585)}
\newcommand{\SensFindOrigHThreeMThreePt}{$-$0.0652}
\newcommand{\SensMacroHOneMFour}{+0.0063~($-$0.0022, +0.0143)}
\newcommand{\SensMacroHOneMFourPt}{+0.0063}
\newcommand{\SensMacroHOneMOne}{+0.0474~(+0.0362, +0.0588)}
\newcommand{\SensMacroHOneMOnePt}{+0.0474}
\newcommand{\SensMacroHOneMThree}{+0.0188~(+0.0089, +0.0287)}
\newcommand{\SensMacroHOneMThreePt}{+0.0188}
\newcommand{\SensMacroHOneMTwo}{+0.0157~(+0.0044, +0.0271)}
\newcommand{\SensMacroHOneMTwoPt}{+0.0157}
\newcommand{\SensMacroHThreeMFour}{$-$0.0412~($-$0.0503, $-$0.0321)}
\newcommand{\SensMacroHThreeMFourPt}{$-$0.0412}
\newcommand{\SensMacroHThreeMThree}{$-$0.0287~($-$0.0370, $-$0.0201)}
\newcommand{\SensMacroHThreeMThreePt}{$-$0.0287}
\newcommand{\SensOrigHOneMFour}{+0.0384~(+0.0336, +0.0430)}
\newcommand{\SensOrigHOneMFourPt}{+0.0384}
\newcommand{\SensOrigHOneMOne}{+0.0734~(+0.0668, +0.0802)}
\newcommand{\SensOrigHOneMOnePt}{+0.0734}
\newcommand{\SensOrigHOneMThree}{+0.0469~(+0.0408, +0.0531)}
\newcommand{\SensOrigHOneMThreePt}{+0.0469}
\newcommand{\SensOrigHOneMTwo}{+0.0763~(+0.0687, +0.0838)}
\newcommand{\SensOrigHOneMTwoPt}{+0.0763}
\newcommand{\SensOrigHThreeMFour}{$-$0.0350~($-$0.0408, $-$0.0292)}
\newcommand{\SensOrigHThreeMFourPt}{$-$0.0350}
\newcommand{\SensOrigHThreeMThree}{$-$0.0265~($-$0.0321, $-$0.0212)}
\newcommand{\SensOrigHThreeMThreePt}{$-$0.0265}
\newcommand{\SensSeedHOneMFour}{$-$0.0469~($-$0.0586, $-$0.0350)}
\newcommand{\SensSeedHOneMFourPt}{$-$0.0469}
\newcommand{\SensSeedHOneMOne}{$-$0.0082~($-$0.0194, +0.0027)}
\newcommand{\SensSeedHOneMOnePt}{$-$0.0082}
\newcommand{\SensSeedHOneMThree}{$-$0.0402~($-$0.0503, $-$0.0298)}
\newcommand{\SensSeedHOneMThreePt}{$-$0.0402}
\newcommand{\SensSeedHOneMTwo}{$-$0.0279~($-$0.0411, $-$0.0147)}
\newcommand{\SensSeedHOneMTwoPt}{$-$0.0279}
\newcommand{\SensSeedHThreeMFour}{$-$0.0386~($-$0.0473, $-$0.0302)}
\newcommand{\SensSeedHThreeMFourPt}{$-$0.0386}
\newcommand{\SensSeedHThreeMThree}{$-$0.0320~($-$0.0405, $-$0.0237)}
\newcommand{\SensSeedHThreeMThreePt}{$-$0.0320}
\newcommand{\SensSeedOrigHOneMFour}{+0.0246~(+0.0182, +0.0315)}
\newcommand{\SensSeedOrigHOneMFourPt}{+0.0246}
\newcommand{\SensSeedOrigHOneMOne}{+0.0511~(+0.0450, +0.0568)}
\newcommand{\SensSeedOrigHOneMOnePt}{+0.0511}
\newcommand{\SensSeedOrigHOneMThree}{+0.0145~(+0.0086, +0.0208)}
\newcommand{\SensSeedOrigHOneMThreePt}{+0.0145}
\newcommand{\SensSeedOrigHOneMTwo}{+0.0669~(+0.0586, +0.0749)}
\newcommand{\SensSeedOrigHOneMTwoPt}{+0.0669}
\newcommand{\SensSeedOrigHThreeMFour}{$-$0.0265~($-$0.0316, $-$0.0216)}
\newcommand{\SensSeedOrigHThreeMFourPt}{$-$0.0265}
\newcommand{\SensSeedOrigHThreeMThree}{$-$0.0366~($-$0.0414, $-$0.0318)}
\newcommand{\SensSeedOrigHThreeMThreePt}{$-$0.0366}
\newcommand{\SensTFiveHOneMFour}{$-$0.0057~($-$0.0146, +0.0033)}
\newcommand{\SensTFiveHOneMFourPt}{$-$0.0057}
\newcommand{\SensTFiveHOneMOne}{+0.0133~(+0.0009, +0.0260)}
\newcommand{\SensTFiveHOneMOnePt}{+0.0133}
\newcommand{\SensTFiveHOneMThree}{$-$0.0114~($-$0.0218, $-$0.0008)}
\newcommand{\SensTFiveHOneMThreePt}{$-$0.0114}
\newcommand{\SensTFiveHOneMTwo}{$-$0.0235~($-$0.0361, $-$0.0111)}
\newcommand{\SensTFiveHOneMTwoPt}{$-$0.0235}
\newcommand{\SensTFiveHThreeMFour}{$-$0.0190~($-$0.0295, $-$0.0082)}
\newcommand{\SensTFiveHThreeMFourPt}{$-$0.0190}
\newcommand{\SensTFiveHThreeMThree}{$-$0.0248~($-$0.0340, $-$0.0151)}
\newcommand{\SensTFiveHThreeMThreePt}{$-$0.0248}
\newcommand{\SensTTenHOneMFour}{$-$0.0009~($-$0.0118, +0.0096)}
\newcommand{\SensTTenHOneMFourPt}{$-$0.0009}
\newcommand{\SensTTenHOneMOne}{+0.0251~(+0.0107, +0.0393)}
\newcommand{\SensTTenHOneMOnePt}{+0.0251}
\newcommand{\SensTTenHOneMThree}{$-$0.0130~($-$0.0269, $-$0.0005)}
\newcommand{\SensTTenHOneMThreePt}{$-$0.0130}
\newcommand{\SensTTenHOneMTwo}{$-$0.0465~($-$0.0632, $-$0.0305)}
\newcommand{\SensTTenHOneMTwoPt}{$-$0.0465}
\newcommand{\SensTTenHThreeMFour}{$-$0.0260~($-$0.0388, $-$0.0120)}
\newcommand{\SensTTenHThreeMFourPt}{$-$0.0260}
\newcommand{\SensTTenHThreeMThree}{$-$0.0381~($-$0.0504, $-$0.0262)}
\newcommand{\SensTTenHThreeMThreePt}{$-$0.0381}
\newcommand{\SensUncNegHOneMFour}{+0.0475~(+0.0388, +0.0563)}
\newcommand{\SensUncNegHOneMFourPt}{+0.0475}
\newcommand{\SensUncNegHOneMOne}{+0.0156~(+0.0051, +0.0270)}
\newcommand{\SensUncNegHOneMOnePt}{+0.0156}
\newcommand{\SensUncNegHOneMThree}{+0.0232~(+0.0141, +0.0331)}
\newcommand{\SensUncNegHOneMThreePt}{+0.0232}
\newcommand{\SensUncNegHOneMTwo}{+0.0184~(+0.0077, +0.0292)}
\newcommand{\SensUncNegHOneMTwoPt}{+0.0184}
\newcommand{\SensUncNegHThreeMFour}{+0.0319~(+0.0219, +0.0417)}
\newcommand{\SensUncNegHThreeMFourPt}{+0.0319}
\newcommand{\SensUncNegHThreeMThree}{+0.0075~($-$0.0014, +0.0163)}
\newcommand{\SensUncNegHThreeMThreePt}{+0.0075}
\newcommand{\SensUncPosHOneMFour}{$-$0.0176~($-$0.0269, $-$0.0091)}
\newcommand{\SensUncPosHOneMFourPt}{$-$0.0176}
\newcommand{\SensUncPosHOneMOne}{+0.0270~(+0.0139, +0.0395)}
\newcommand{\SensUncPosHOneMOnePt}{+0.0270}
\newcommand{\SensUncPosHOneMThree}{$-$0.0025~($-$0.0135, +0.0080)}
\newcommand{\SensUncPosHOneMThreePt}{$-$0.0025}
\newcommand{\SensUncPosHOneMTwo}{$-$0.0207~($-$0.0322, $-$0.0092)}
\newcommand{\SensUncPosHOneMTwoPt}{$-$0.0207}
\newcommand{\SensUncPosHThreeMFour}{$-$0.0446~($-$0.0551, $-$0.0339)}
\newcommand{\SensUncPosHThreeMFourPt}{$-$0.0446}
\newcommand{\SensUncPosHThreeMThree}{$-$0.0295~($-$0.0384, $-$0.0205)}
\newcommand{\SensUncPosHThreeMThreePt}{$-$0.0295}
\newcommand{\SensUnmNegHOneMFour}{+0.1306~(+0.1211, +0.1406)}
\newcommand{\SensUnmNegHOneMFourPt}{+0.1306}
\newcommand{\SensUnmNegHOneMOne}{+0.0649~(+0.0532, +0.0774)}
\newcommand{\SensUnmNegHOneMOnePt}{+0.0649}
\newcommand{\SensUnmNegHOneMThree}{+0.0828~(+0.0724, +0.0942)}
\newcommand{\SensUnmNegHOneMThreePt}{+0.0828}
\newcommand{\SensUnmNegHOneMTwo}{+0.0729~(+0.0665, +0.0796)}
\newcommand{\SensUnmNegHOneMTwoPt}{+0.0729}
\newcommand{\SensUnmNegHThreeMFour}{+0.0657~(+0.0588, +0.0725)}
\newcommand{\SensUnmNegHThreeMFourPt}{+0.0657}
\newcommand{\SensUnmNegHThreeMThree}{+0.0179~(+0.0127, +0.0232)}
\newcommand{\SensUnmNegHThreeMThreePt}{+0.0179}
\newcommand{\SrcAtelPosRate}{97.7\%}
\newcommand{\SrcEdemaUnmNeg}{13.8\%}
\newcommand{\SrcEffUnmNeg}{22.9\%}
\newcommand{\StudiesExtFindings}{132,923}
\newcommand{\StudiesExtImpression}{132,923}
\newcommand{\StudiesSrcFindings}{9,293}
\newcommand{\StudiesSrcImpression}{14,766}
\newcommand{\ThrSDMaxMFour}{0.077}
\newcommand{\ViewAPMFour}{$-$0.0075~($-$0.0178, +0.0020)}
\newcommand{\ViewPAMFour}{+0.0335~(+0.0169, +0.0487)}

\renewcommand{\thetable}{S\arabic{table}}
\renewcommand{\thesection}{S\arabic{section}}
\makeatletter\renewcommand*\l@section{\@dottedtocline{1}{0pt}{2.6em}}\makeatother

\title{Supplementary methods and results\\[2pt]\large Frozen selective prediction
for multimodal chest-radiograph models across two hospitals (Revision R1)}
\author{}
\date{6 October 2026}

\begin{document}
\maketitle
\tableofcontents

\section{Estimators and their implementation}
Notation follows the main text: study $i$, pathology $j$, mean frontal-image
probability $p_{ij}$, frozen threshold $t_j$, decision $\hat y_{ij}$, label
$y_{ij}$, mask $m_{ij}$ (1 for positive or negative; 0 for uncertain or
unmentioned), $n_i=\sum_jm_{ij}$, $w_i=\sum_jm_{ij}\mathbb{1}[\hat y_{ij}\neq y_{ij}]$,
acceptance $a_i=\mathbb{1}[\min_j|2p_{ij}-1|\ge\tau]$.
\begin{center}\small
\begin{tabular}{p{0.2\linewidth}p{0.42\linewidth}p{0.3\linewidth}}
\toprule
Estimator & Definition & Code\\
\midrule
original & $\sum_i a_i w_i/\max(n_i,1)\,/\,\sum_i a_i$ & \texttt{hamming\_error} (calibration.policy) with \texttt{selective\_error} (analysis.bootstrap)\\
evaluable-study (corrected) & $\sum_i a_i\mathbb{1}[n_i>0]w_i/n_i\,/\,\sum_i a_i\mathbb{1}[n_i>0]$ & \texttt{revision.estimators}\\
cell-pooled & $\sum_i a_iw_i/\sum_ia_in_i$ & \texttt{revision.estimators}\\
pathology-macro & $\frac15\sum_j\frac{\sum_ia_im_{ij}\mathbb{1}[\hat y_{ij}\ne y_{ij}]}{\sum_ia_im_{ij}}$ & \texttt{revision.estimators}\\
coverage (all) & $\frac1N\sum_ia_i$ & \\
coverage (evaluable) & $\sum_ia_i\mathbb{1}[n_i>0]/\sum_i\mathbb{1}[n_i>0]$ & \\
\bottomrule
\end{tabular}
\end{center}
Label variants change only $m$ and $y$: \emph{unmentioned$\to$negative} sets blank
cells to observed negatives; \emph{uncertain$\to$positive} or
\emph{$\to$negative} sets $-1$ cells accordingly. Uncertain-label variants were
SAP-registered sensitivities; the unmentioned variant is a post hoc bracket,
not a reference standard.

Bootstrap replicates are computed as patient-multiplicity-weighted sums over
patient-aggregated term matrices. The random integers are drawn in exactly the
order of the Stage 10 code, so every original interval is reproduced
(\cref{tab:repro}). Unit tests pin the equivalence on synthetic data.

\section{Reproduction and alignment checks}
The Stage 8/9b caches do not store study identifiers. Study identity was
recovered by replaying cohort construction through the original functions and
verified row by row against every cache: patient order, observed labels and
masks must match exactly. At the external site, one image present at the time
of the primary run has since been lost from storage. Its inclusion was required
to reproduce the primary caches, and its absence reproduces the seed-replicate
caches (132,922 studies).
\begin{table}[h]\centering\small
\caption{Reproduction checks.}\label{tab:repro}
\begin{tabular}{llr}
\toprule
Check & Item & Result\\
\midrule
Stage 10 replay & primary (14 rows) & max $|\Delta|$ = 0\\
Stage 10 replay & T5 (14 rows) & max $|\Delta|$ = 0\\
Stage 10 replay & T10 (14 rows) & max $|\Delta|$ = 0\\
Stage 10 replay & seed20260719 (7 rows) & max $|\Delta|$ = 0\\
Stage 10 replay & findings S1 (14 rows) & max $|\Delta|$ = 0\\
Stage 7 re-selection & M1 & thresholds $|\Delta|$=0; cutoffs $|\Delta|$=0\\
Stage 7 re-selection & M2 & thresholds $|\Delta|$=0; cutoffs $|\Delta|$=0\\
Stage 7 re-selection & M3 & thresholds $|\Delta|$=0; cutoffs $|\Delta|$=0\\
Stage 7 re-selection & M4 & thresholds $|\Delta|$=0; cutoffs $|\Delta|$=0\\
Stage 7 re-selection & M1 seed 20260719 & thresholds $|\Delta|$=0; cutoffs $|\Delta|$=0\\
Stage 7 re-selection & M4 seed 20260719 & thresholds $|\Delta|$=0; cutoffs $|\Delta|$=0\\
Cohort alignment & source impression $T=3$ & 14,766 rows exact\\
Cohort alignment & source findings $T=3$ & 9,293 rows exact\\
Cohort alignment & external impression $T=3$ & 132,923 rows exact\\
Cohort alignment & external findings $T=3$ & 132,923 rows exact\\
Cohort alignment & source impression $T=5$ & 13,600 rows exact\\
Cohort alignment & source impression $T=10$ & 8,381 rows exact\\
Cohort alignment & external impression $T=5$ & 111,040 rows exact\\
Cohort alignment & external impression $T=10$ & 23,146 rows exact\\
Cohort alignment & source impression $T=3$ seed 20260719 & 14,766 rows exact\\
Cohort alignment & external impression $T=3$ seed 20260719 & 132,922 rows exact\\
\bottomrule
\end{tabular}

\end{table}

\begin{landscape}
\section{Every contrast under every specification}
\footnotesize
\begin{longtable}{llrrrr}
\toprule
Contrast & Specification & M1 & M2 & M3 & M4\\
\midrule
\endhead
H1 & original & +0.073 (+0.067, +0.080) & +0.076 (+0.069, +0.084) & +0.047 (+0.041, +0.053) & +0.038 (+0.034, +0.043)\\
 & evaluable-study & +0.012 ($-$0.000, +0.024) & $-$0.017 ($-$0.029, $-$0.005) & $-$0.004 ($-$0.014, +0.006) & $-$0.002 ($-$0.011, +0.006)\\
 & cell-pooled & +0.016 (+0.005, +0.027) & $-$0.001 ($-$0.014, +0.011) & $-$0.003 ($-$0.012, +0.007) & +0.008 ($-$0.000, +0.015)\\
 & pathology-macro & +0.047 (+0.036, +0.059) & +0.016 (+0.004, +0.027) & +0.019 (+0.009, +0.029) & +0.006 ($-$0.002, +0.014)\\
 & unc.$\to$neg & +0.016 (+0.005, +0.027) & +0.018 (+0.008, +0.029) & +0.023 (+0.014, +0.033) & +0.047 (+0.039, +0.056)\\
 & unc.$\to$pos & +0.027 (+0.014, +0.039) & $-$0.021 ($-$0.032, $-$0.009) & $-$0.003 ($-$0.014, +0.008) & $-$0.018 ($-$0.027, $-$0.009)\\
 & unm.$\to$neg & +0.065 (+0.053, +0.077) & +0.073 (+0.066, +0.080) & +0.083 (+0.072, +0.094) & +0.131 (+0.121, +0.141)\\
 & cov.\ 70\% & +0.008 ($-$0.005, +0.021) & $-$0.012 ($-$0.025, +0.000) & $-$0.007 ($-$0.017, +0.003) & +0.003 ($-$0.006, +0.011)\\
 & cov.\ 90\% & +0.017 (+0.005, +0.029) & $-$0.018 ($-$0.030, $-$0.007) & $-$0.001 ($-$0.011, +0.009) & $-$0.006 ($-$0.015, +0.002)\\
 & $T=5$ & +0.013 (+0.001, +0.026) & $-$0.024 ($-$0.036, $-$0.011) & $-$0.011 ($-$0.022, $-$0.001) & $-$0.006 ($-$0.015, +0.003)\\
 & $T=10$ & +0.025 (+0.011, +0.039) & $-$0.046 ($-$0.063, $-$0.031) & $-$0.013 ($-$0.027, $-$0.000) & $-$0.001 ($-$0.012, +0.010)\\
 & findings & +0.095 (+0.087, +0.103) & $-$0.000 ($-$0.011, +0.010) & +0.029 (+0.019, +0.038) & +0.005 ($-$0.003, +0.013)\\
 & findings, original & $-$0.067 ($-$0.074, $-$0.061) & $-$0.238 ($-$0.247, $-$0.230) & $-$0.132 ($-$0.140, $-$0.125) & $-$0.108 ($-$0.114, $-$0.101)\\
 & seed 2 & $-$0.008 ($-$0.019, +0.003) & $-$0.028 ($-$0.041, $-$0.015) & $-$0.040 ($-$0.050, $-$0.030) & $-$0.047 ($-$0.059, $-$0.035)\\
 & seed 2, original & +0.051 (+0.045, +0.057) & +0.067 (+0.059, +0.075) & +0.015 (+0.009, +0.021) & +0.025 (+0.018, +0.031)\\
\midrule
H2-C1 & original & -- & -- & +0.002 ($-$0.003, +0.008) & +0.010 (+0.005, +0.014)\\
 & evaluable-study & -- & -- & +0.005 ($-$0.004, +0.013) & +0.009 (+0.000, +0.017)\\
 & cell-pooled & -- & -- & +0.004 ($-$0.003, +0.012) & $-$0.007 ($-$0.014, +0.000)\\
 & pathology-macro & -- & -- & +0.006 ($-$0.003, +0.015) & $-$0.020 ($-$0.028, $-$0.012)\\
 & unc.$\to$neg & -- & -- & $-$0.022 ($-$0.030, $-$0.014) & $-$0.018 ($-$0.026, $-$0.009)\\
 & unc.$\to$pos & -- & -- & +0.026 (+0.018, +0.034) & +0.028 (+0.020, +0.036)\\
 & unm.$\to$neg & -- & -- & $-$0.033 ($-$0.037, $-$0.028) & $-$0.035 ($-$0.040, $-$0.030)\\
 & cov.\ 70\% & -- & -- & +0.004 ($-$0.006, +0.013) & +0.002 ($-$0.008, +0.011)\\
 & cov.\ 90\% & -- & -- & +0.005 ($-$0.003, +0.013) & +0.012 (+0.004, +0.019)\\
 & $T=5$ & -- & -- & +0.013 (+0.004, +0.022) & +0.012 (+0.003, +0.021)\\
 & $T=10$ & -- & -- & +0.018 (+0.006, +0.028) & +0.005 ($-$0.007, +0.015)\\
 & findings & -- & -- & +0.035 (+0.027, +0.042) & +0.111 (+0.103, +0.118)\\
 & findings, original & -- & -- & +0.034 (+0.028, +0.040) & +0.062 (+0.056, +0.069)\\
 & seed 2 & -- & -- & +0.015 (+0.007, +0.023) & +0.003 ($-$0.005, +0.011)\\
 & seed 2, original & -- & -- & +0.015 (+0.011, +0.020) & +0.009 (+0.004, +0.014)\\
\midrule
H2-C2 & original & -- & -- & $-$0.022 ($-$0.028, $-$0.016) & $-$0.014 ($-$0.019, $-$0.009)\\
 & evaluable-study & -- & -- & $-$0.047 ($-$0.057, $-$0.037) & $-$0.029 ($-$0.038, $-$0.020)\\
 & cell-pooled & -- & -- & $-$0.043 ($-$0.053, $-$0.035) & $-$0.025 ($-$0.033, $-$0.017)\\
 & pathology-macro & -- & -- & $-$0.045 ($-$0.056, $-$0.036) & $-$0.027 ($-$0.035, $-$0.018)\\
 & unc.$\to$neg & -- & -- & $-$0.039 ($-$0.049, $-$0.030) & $-$0.030 ($-$0.038, $-$0.021)\\
 & unc.$\to$pos & -- & -- & $-$0.045 ($-$0.054, $-$0.035) & $-$0.023 ($-$0.031, $-$0.014)\\
 & unm.$\to$neg & -- & -- & +0.010 (+0.005, +0.016) & +0.016 (+0.010, +0.021)\\
 & cov.\ 70\% & -- & -- & $-$0.049 ($-$0.060, $-$0.038) & $-$0.033 ($-$0.042, $-$0.023)\\
 & cov.\ 90\% & -- & -- & $-$0.050 ($-$0.059, $-$0.040) & $-$0.029 ($-$0.037, $-$0.021)\\
 & $T=5$ & -- & -- & $-$0.048 ($-$0.058, $-$0.037) & $-$0.028 ($-$0.036, $-$0.019)\\
 & $T=10$ & -- & -- & $-$0.037 ($-$0.052, $-$0.024) & $-$0.029 ($-$0.040, $-$0.018)\\
 & findings & -- & -- & +0.000 ($-$0.008, +0.009) & +0.015 (+0.007, +0.023)\\
 & findings, original & -- & -- & $-$0.018 ($-$0.025, $-$0.011) & +0.004 ($-$0.002, +0.011)\\
 & seed 2 & -- & -- & $-$0.040 ($-$0.049, $-$0.032) & $-$0.013 ($-$0.021, $-$0.005)\\
 & seed 2, original & -- & -- & $-$0.018 ($-$0.023, $-$0.012) & $-$0.006 ($-$0.011, $-$0.001)\\
\midrule
H3 & original & -- & -- & $-$0.027 ($-$0.032, $-$0.021) & $-$0.035 ($-$0.041, $-$0.029)\\
 & evaluable-study & -- & -- & $-$0.016 ($-$0.026, $-$0.007) & $-$0.014 ($-$0.024, $-$0.003)\\
 & cell-pooled & -- & -- & $-$0.018 ($-$0.026, $-$0.010) & $-$0.008 ($-$0.017, +0.001)\\
 & pathology-macro & -- & -- & $-$0.029 ($-$0.037, $-$0.020) & $-$0.041 ($-$0.050, $-$0.032)\\
 & unc.$\to$neg & -- & -- & +0.008 ($-$0.001, +0.016) & +0.032 (+0.022, +0.042)\\
 & unc.$\to$pos & -- & -- & $-$0.029 ($-$0.038, $-$0.020) & $-$0.045 ($-$0.055, $-$0.034)\\
 & unm.$\to$neg & -- & -- & +0.018 (+0.013, +0.023) & +0.066 (+0.059, +0.073)\\
 & cov.\ 70\% & -- & -- & $-$0.015 ($-$0.025, $-$0.005) & $-$0.005 ($-$0.016, +0.006)\\
 & cov.\ 90\% & -- & -- & $-$0.018 ($-$0.027, $-$0.009) & $-$0.023 ($-$0.033, $-$0.012)\\
 & $T=5$ & -- & -- & $-$0.025 ($-$0.034, $-$0.015) & $-$0.019 ($-$0.030, $-$0.008)\\
 & $T=10$ & -- & -- & $-$0.038 ($-$0.050, $-$0.026) & $-$0.026 ($-$0.039, $-$0.012)\\
 & findings & -- & -- & $-$0.067 ($-$0.074, $-$0.058) & $-$0.090 ($-$0.098, $-$0.082)\\
 & findings, original & -- & -- & $-$0.065 ($-$0.071, $-$0.058) & $-$0.041 ($-$0.047, $-$0.034)\\
 & seed 2 & -- & -- & $-$0.032 ($-$0.040, $-$0.024) & $-$0.039 ($-$0.047, $-$0.030)\\
 & seed 2, original & -- & -- & $-$0.037 ($-$0.041, $-$0.032) & $-$0.026 ($-$0.032, $-$0.022)\\
\midrule
H4 source & original & -- & -- & +0.047 (+0.042, +0.052) & +0.017 (+0.013, +0.022)\\
 & evaluable-study & -- & -- & +0.078 (+0.069, +0.087) & +0.030 (+0.021, +0.038)\\
 & cell-pooled & -- & -- & +0.077 (+0.069, +0.085) & +0.022 (+0.015, +0.030)\\
 & pathology-macro & -- & -- & +0.076 (+0.067, +0.085) & +0.019 (+0.011, +0.028)\\
 & unc.$\to$neg & -- & -- & +0.054 (+0.045, +0.062) & +0.024 (+0.016, +0.033)\\
 & unc.$\to$pos & -- & -- & +0.083 (+0.074, +0.092) & +0.032 (+0.024, +0.040)\\
 & unm.$\to$neg & -- & -- & +0.018 (+0.013, +0.023) & +0.022 (+0.017, +0.026)\\
 & cov.\ 70\% & -- & -- & +0.078 (+0.069, +0.088) & +0.028 (+0.020, +0.037)\\
 & cov.\ 90\% & -- & -- & +0.081 (+0.073, +0.089) & +0.033 (+0.025, +0.040)\\
 & $T=5$ & -- & -- & +0.083 (+0.073, +0.092) & +0.034 (+0.026, +0.043)\\
 & $T=10$ & -- & -- & +0.088 (+0.076, +0.099) & +0.031 (+0.021, +0.042)\\
 & findings & -- & -- & +0.065 (+0.058, +0.072) & +0.052 (+0.045, +0.058)\\
 & findings, original & -- & -- & +0.061 (+0.054, +0.067) & +0.048 (+0.042, +0.054)\\
 & seed 2 & -- & -- & +0.056 (+0.048, +0.065) & +0.004 ($-$0.003, +0.012)\\
 & seed 2, original & -- & -- & +0.036 (+0.031, +0.040) & +0.006 (+0.002, +0.010)\\
\midrule
H4 external & original & -- & -- & +0.023 (+0.021, +0.025) & $-$0.006 ($-$0.008, $-$0.004)\\
 & evaluable-study & -- & -- & +0.026 (+0.024, +0.029) & $-$0.008 ($-$0.010, $-$0.005)\\
 & cell-pooled & -- & -- & +0.029 (+0.027, +0.031) & +0.005 (+0.002, +0.007)\\
 & pathology-macro & -- & -- & +0.025 (+0.022, +0.027) & +0.013 (+0.011, +0.015)\\
 & unc.$\to$neg & -- & -- & +0.037 (+0.034, +0.039) & +0.013 (+0.010, +0.015)\\
 & unc.$\to$pos & -- & -- & +0.013 (+0.011, +0.015) & $-$0.019 ($-$0.021, $-$0.017)\\
 & unm.$\to$neg & -- & -- & +0.061 (+0.060, +0.063) & +0.073 (+0.071, +0.074)\\
 & cov.\ 70\% & -- & -- & +0.026 (+0.023, +0.029) & $-$0.006 ($-$0.009, $-$0.004)\\
 & cov.\ 90\% & -- & -- & +0.027 (+0.025, +0.028) & $-$0.008 ($-$0.010, $-$0.006)\\
 & $T=5$ & -- & -- & +0.022 (+0.020, +0.025) & $-$0.006 ($-$0.008, $-$0.003)\\
 & $T=10$ & -- & -- & +0.033 (+0.027, +0.039) & $-$0.003 ($-$0.009, +0.003)\\
 & findings & -- & -- & +0.031 (+0.027, +0.035) & $-$0.044 ($-$0.049, $-$0.040)\\
 & findings, original & -- & -- & +0.009 (+0.008, +0.010) & $-$0.010 ($-$0.012, $-$0.009)\\
 & seed 2 & -- & -- & +0.001 ($-$0.001, +0.003) & $-$0.011 ($-$0.013, $-$0.009)\\
 & seed 2, original & -- & -- & +0.002 (+0.001, +0.004) & $-$0.009 ($-$0.011, $-$0.007)\\
\bottomrule
\end{longtable}

\normalsize
\end{landscape}

\section{Pathology-specific contrasts and multiplicity}
The original SAP required pathology-specific reporting with Holm correction
across the five pathologies within each contrast family. Neither was performed
before the external results were seen. Both are reported here post hoc.
\begin{table}[h]\centering\small
\caption{Pathology-specific H1 (selective error over labelled cells of that
pathology among accepted studies). Holm $p$: bootstrap $p$-values adjusted
across the five pathologies.}
\begin{adjustbox}{max width=\linewidth}\begin{tabular}{llrrr}
\toprule
Model & Pathology & Primary labels (95\% CI) & Holm $p$ & Unmentioned $\to$ negative\\
\midrule
M3 & Cardiomegaly & $-$0.018 ($-$0.034, $-$0.001) & 0.080 & +0.103 (+0.088, +0.118)\\
 & Edema & +0.031 (+0.014, +0.048) & 0.012 & +0.083 (+0.070, +0.098)\\
 & Consolidation & $-$0.030 ($-$0.054, $-$0.005) & 0.057 & +0.148 (+0.128, +0.171)\\
 & Atelectasis & +0.114 (+0.087, +0.142) & 0.005 & +0.103 (+0.088, +0.119)\\
 & Pleural Effusion & $-$0.004 ($-$0.016, +0.008) & 0.530 & $-$0.008 ($-$0.020, +0.005)\\
\addlinespace
M4 & Cardiomegaly & +0.053 (+0.036, +0.068) & 0.005 & +0.179 (+0.165, +0.194)\\
 & Edema & +0.018 (+0.000, +0.033) & 0.141 & +0.181 (+0.167, +0.196)\\
 & Consolidation & $-$0.013 ($-$0.036, +0.008) & 0.492 & +0.129 (+0.110, +0.150)\\
 & Atelectasis & $-$0.026 ($-$0.045, $-$0.010) & 0.008 & +0.153 (+0.138, +0.166)\\
 & Pleural Effusion & +0.000 ($-$0.012, +0.013) & 0.933 & +0.030 (+0.019, +0.042)\\
\bottomrule
\end{tabular}
\end{adjustbox}
\end{table}
\begin{table}[h]\centering\small
\caption{Multiplicity sensitivity for the six registered confirmatory contrasts
(post hoc).}
\begin{adjustbox}{max width=\linewidth}\begin{tabular}{lllrrr}
\toprule
Estimator & H & Model & Estimate & 95\% CI & Bonferroni ($k=6$) 99.17\% CI\\
\midrule
original & H1 & M3 & +0.0469 & (+0.0408, +0.0531) & (+0.0388, +0.0555)\\
original & H1 & M4 & +0.0384 & (+0.0336, +0.0430) & (+0.0315, +0.0444)\\
original & H2-C1 & M3 & +0.0024 & ($-$0.0027, +0.0076) & ($-$0.0046, +0.0100)\\
original & H2-C1 & M4 & +0.0096 & (+0.0048, +0.0144) & (+0.0033, +0.0158)\\
original & H3 & M3 & $-$0.0265 & ($-$0.0321, $-$0.0212) & ($-$0.0340, $-$0.0194)\\
original & H3 & M4 & $-$0.0350 & ($-$0.0408, $-$0.0292) & ($-$0.0429, $-$0.0276)\\
\addlinespace
evaluable-study & H1 & M3 & $-$0.0042 & ($-$0.0145, +0.0061) & ($-$0.0182, +0.0100)\\
evaluable-study & H1 & M4 & $-$0.0018 & ($-$0.0108, +0.0064) & ($-$0.0141, +0.0090)\\
evaluable-study & H2-C1 & M3 & +0.0048 & ($-$0.0038, +0.0134) & ($-$0.0072, +0.0178)\\
evaluable-study & H2-C1 & M4 & +0.0089 & (+0.0002, +0.0173) & ($-$0.0018, +0.0198)\\
evaluable-study & H3 & M3 & $-$0.0162 & ($-$0.0256, $-$0.0070) & ($-$0.0285, $-$0.0035)\\
evaluable-study & H3 & M4 & $-$0.0138 & ($-$0.0243, $-$0.0029) & ($-$0.0277, +0.0005)\\
\bottomrule
\end{tabular}
\end{adjustbox}
\end{table}

\begin{landscape}
\section{Controlled label-source decomposition (all models, with intervals)}
Step codes as in the main text: A0 primary; A1 findings labels only; A1c
common labelled cells with impression (A1c-I) or findings (A1c-F) labels; A2x findings cohort
only; A2 findings cohort and labels; A3 plus findings thresholds (published
S1); B2 findings labels and thresholds on the primary cohort.
\footnotesize
\begin{longtable}{llrrrr}
\toprule
Step & Estimator & M1 & M2 & M3 & M4\\
\midrule
\endhead
A0 & original & +0.073 (+0.067, +0.080) & +0.076 (+0.069, +0.084) & +0.047 (+0.041, +0.053) & +0.038 (+0.034, +0.043)\\
 & evaluable-study & +0.012 ($-$0.000, +0.024) & $-$0.017 ($-$0.029, $-$0.005) & $-$0.004 ($-$0.014, +0.006) & $-$0.002 ($-$0.011, +0.006)\\
 & cell-micro & +0.016 (+0.005, +0.027) & $-$0.001 ($-$0.014, +0.011) & $-$0.003 ($-$0.012, +0.007) & +0.008 ($-$0.000, +0.015)\\
\midrule
A1 & original & $-$0.027 ($-$0.031, $-$0.022) & $-$0.241 ($-$0.253, $-$0.230) & $-$0.061 ($-$0.067, $-$0.056) & $-$0.084 ($-$0.090, $-$0.078)\\
 & evaluable-study & +0.100 (+0.091, +0.108) & $-$0.213 ($-$0.225, $-$0.201) & +0.009 ($-$0.000, +0.018) & $-$0.039 ($-$0.048, $-$0.030)\\
 & cell-micro & +0.086 (+0.078, +0.094) & $-$0.184 ($-$0.196, $-$0.172) & +0.006 ($-$0.003, +0.015) & $-$0.031 ($-$0.040, $-$0.023)\\
\midrule
A1c-I & original & $-$0.002 ($-$0.006, +0.002) & $-$0.020 ($-$0.024, $-$0.016) & $-$0.006 ($-$0.010, $-$0.002) & $-$0.006 ($-$0.009, $-$0.004)\\
 & evaluable-study & +0.020 ($-$0.002, +0.040) & $-$0.072 ($-$0.089, $-$0.056) & $-$0.004 ($-$0.021, +0.013) & $-$0.021 ($-$0.035, $-$0.009)\\
 & cell-micro & +0.032 (+0.012, +0.051) & $-$0.060 ($-$0.077, $-$0.044) & +0.006 ($-$0.009, +0.022) & $-$0.015 ($-$0.028, $-$0.004)\\
\midrule
A1c-F & original & $-$0.003 ($-$0.007, +0.001) & $-$0.021 ($-$0.025, $-$0.017) & $-$0.007 ($-$0.011, $-$0.003) & $-$0.008 ($-$0.010, $-$0.005)\\
 & evaluable-study & +0.016 ($-$0.006, +0.036) & $-$0.076 ($-$0.093, $-$0.058) & $-$0.009 ($-$0.026, +0.009) & $-$0.027 ($-$0.040, $-$0.014)\\
 & cell-micro & +0.028 (+0.009, +0.047) & $-$0.063 ($-$0.080, $-$0.047) & +0.002 ($-$0.013, +0.019) & $-$0.020 ($-$0.032, $-$0.008)\\
\midrule
A2x & original & +0.098 (+0.092, +0.105) & +0.145 (+0.138, +0.151) & +0.081 (+0.074, +0.088) & +0.066 (+0.061, +0.071)\\
 & evaluable-study & $-$0.014 ($-$0.031, +0.003) & +0.045 (+0.029, +0.060) & +0.001 ($-$0.014, +0.016) & +0.013 ($-$0.000, +0.024)\\
 & cell-micro & $-$0.006 ($-$0.021, +0.010) & +0.076 (+0.061, +0.091) & +0.014 ($-$0.000, +0.027) & +0.025 (+0.013, +0.036)\\
\midrule
A2 & original & $-$0.075 ($-$0.082, $-$0.068) & $-$0.409 ($-$0.419, $-$0.398) & $-$0.136 ($-$0.143, $-$0.128) & $-$0.161 ($-$0.168, $-$0.153)\\
 & evaluable-study & +0.100 (+0.091, +0.108) & $-$0.213 ($-$0.225, $-$0.200) & +0.009 ($-$0.000, +0.019) & $-$0.039 ($-$0.048, $-$0.030)\\
 & cell-micro & +0.086 (+0.077, +0.093) & $-$0.184 ($-$0.195, $-$0.173) & +0.006 ($-$0.003, +0.015) & $-$0.031 ($-$0.039, $-$0.023)\\
\midrule
A3 & original & $-$0.067 ($-$0.074, $-$0.061) & $-$0.238 ($-$0.247, $-$0.230) & $-$0.132 ($-$0.140, $-$0.125) & $-$0.108 ($-$0.114, $-$0.101)\\
 & evaluable-study & +0.095 (+0.087, +0.103) & $-$0.000 ($-$0.011, +0.010) & +0.029 (+0.019, +0.038) & +0.005 ($-$0.003, +0.013)\\
 & cell-micro & +0.080 (+0.072, +0.088) & +0.008 ($-$0.002, +0.017) & +0.025 (+0.017, +0.034) & +0.001 ($-$0.006, +0.008)\\
\midrule
B2 & original & $-$0.023 ($-$0.028, $-$0.019) & $-$0.128 ($-$0.135, $-$0.120) & $-$0.058 ($-$0.063, $-$0.053) & $-$0.052 ($-$0.057, $-$0.048)\\
 & evaluable-study & +0.095 (+0.087, +0.103) & $-$0.000 ($-$0.011, +0.011) & +0.029 (+0.019, +0.038) & +0.005 ($-$0.003, +0.013)\\
 & cell-micro & +0.080 (+0.073, +0.088) & +0.008 ($-$0.002, +0.018) & +0.025 (+0.016, +0.034) & +0.001 ($-$0.007, +0.009)\\
\bottomrule
\end{longtable}

\normalsize
\end{landscape}

\section{Operating characteristics with unmentioned labels counted as negative}
\begin{table}[h]\centering\small
\caption{Sensitivity and specificity among accepted studies (C0) when unmentioned
cells count as negative.}
\begin{tabular}{llrrrr}
\toprule
Model & Pathology & Sens.\\ src & Spec.\\ src & Sens.\\ ext & Spec.\\ ext\\
\midrule
M3 & Cardiomegaly & 0.872 & 0.473 & 0.965 & 0.345\\
 & Edema & 0.840 & 0.718 & 0.808 & 0.593\\
 & Consolidation & 0.824 & 0.610 & 0.808 & 0.443\\
 & Atelectasis & 0.667 & 0.619 & 0.550 & 0.520\\
 & Pleural Effusion & 0.863 & 0.703 & 0.851 & 0.658\\
M4 & Cardiomegaly & 0.887 & 0.497 & 0.972 & 0.281\\
 & Edema & 0.848 & 0.736 & 0.885 & 0.455\\
 & Consolidation & 0.841 & 0.650 & 0.805 & 0.507\\
 & Atelectasis & 0.886 & 0.326 & 0.911 & 0.150\\
 & Pleural Effusion & 0.870 & 0.636 & 0.950 & 0.419\\
\bottomrule
\end{tabular}

\end{table}

\section{Selective-score diagnostics and comparators}
\begin{table}[h]\centering\small
\caption{AURC and AUGRC over evaluable studies (registered MET008 and MET007,
not computed in the original analysis), external minus source with joint
bootstrap intervals.}
\begin{adjustbox}{max width=\linewidth}\begin{tabular}{lllrrr}
\toprule
Score & Metric & Model & Source & External & Difference (95\\% CI)\\
\midrule
frozen & AURC & M1 & 0.1631 & 0.1750 & +0.0120 ($-$0.0011, +0.0240)\\
frozen & AUGRC & M1 & 0.0961 & 0.1023 & +0.0062 (+0.0001, +0.0122)\\
frozen & AURC & M2 & 0.2373 & 0.2209 & $-$0.0164 ($-$0.0282, $-$0.0044)\\
frozen & AUGRC & M2 & 0.1388 & 0.1277 & $-$0.0111 ($-$0.0171, $-$0.0052)\\
frozen & AURC & M3 & 0.1317 & 0.1269 & $-$0.0048 ($-$0.0142, +0.0048)\\
frozen & AUGRC & M3 & 0.0812 & 0.0813 & +0.0001 ($-$0.0049, +0.0050)\\
frozen & AURC & M4 & 0.1011 & 0.1015 & +0.0005 ($-$0.0078, +0.0083)\\
frozen & AUGRC & M4 & 0.0614 & 0.0613 & $-$0.0001 ($-$0.0043, +0.0039)\\
threshold margin & AURC & M1 & 0.1888 & 0.2153 & +0.0265 (+0.0125, +0.0398)\\
threshold margin & AUGRC & M1 & 0.1035 & 0.1179 & +0.0144 (+0.0078, +0.0206)\\
threshold margin & AURC & M2 & 0.2934 & 0.2664 & $-$0.0269 ($-$0.0401, $-$0.0137)\\
threshold margin & AUGRC & M2 & 0.1555 & 0.1413 & $-$0.0142 ($-$0.0201, $-$0.0080)\\
threshold margin & AURC & M3 & 0.1553 & 0.1523 & $-$0.0030 ($-$0.0146, +0.0080)\\
threshold margin & AUGRC & M3 & 0.0893 & 0.0885 & $-$0.0009 ($-$0.0062, +0.0046)\\
threshold margin & AURC & M4 & 0.1030 & 0.0945 & $-$0.0085 ($-$0.0173, $-$0.0003)\\
threshold margin & AUGRC & M4 & 0.0622 & 0.0578 & $-$0.0044 ($-$0.0087, $-$0.0002)\\
\bottomrule
\end{tabular}
\end{adjustbox}
\end{table}
\begin{table}[h]\centering\small
\caption{Per-cell failure-detection AUROC: does the score rank the correctness
of the thresholded decision for labelled cells? $|2p-1|$ is the frozen score
(centred at 0.5); the logit margin is centred at the decision threshold.}
\begin{adjustbox}{max width=\linewidth}\begin{tabular}{llrrrr}
\toprule
 & & \multicolumn{2}{c}{Source} & \multicolumn{2}{c}{External}\\
\cmidrule(lr){3-4}\cmidrule(lr){5-6}
Model & Pathology & $|2p-1|$ & $|\mathrm{logit}\,p-\mathrm{logit}\,t|$ & $|2p-1|$ & $|\mathrm{logit}\,p-\mathrm{logit}\,t|$\\
\midrule
M1 & Atelectasis & 0.970 & 0.494 & 0.953 & 0.336\\
M1 & Cardiomegaly & 0.802 & 0.806 & 0.807 & 0.840\\
M1 & Consolidation & 0.729 & 0.746 & 0.796 & 0.781\\
M1 & Edema & 0.782 & 0.781 & 0.707 & 0.680\\
M1 & Pleural Effusion & 0.865 & 0.803 & 0.917 & 0.771\\
M2 & Atelectasis & 0.968 & 0.396 & 0.980 & 0.418\\
M2 & Cardiomegaly & 0.748 & 0.687 & 0.590 & 0.570\\
M2 & Consolidation & 0.728 & 0.716 & 0.468 & 0.628\\
M2 & Edema & 0.685 & 0.659 & 0.654 & 0.624\\
M2 & Pleural Effusion & 0.735 & 0.686 & 0.671 & 0.644\\
M3 & Atelectasis & 0.975 & 0.540 & 0.969 & 0.476\\
M3 & Cardiomegaly & 0.786 & 0.786 & 0.670 & 0.739\\
M3 & Consolidation & 0.806 & 0.803 & 0.709 & 0.738\\
M3 & Edema & 0.773 & 0.763 & 0.788 & 0.732\\
M3 & Pleural Effusion & 0.831 & 0.809 & 0.781 & 0.823\\
M4 & Atelectasis & 0.964 & 0.825 & 0.968 & 0.832\\
M4 & Cardiomegaly & 0.842 & 0.822 & 0.774 & 0.808\\
M4 & Consolidation & 0.783 & 0.765 & 0.788 & 0.794\\
M4 & Edema & 0.795 & 0.807 & 0.746 & 0.763\\
M4 & Pleural Effusion & 0.872 & 0.801 & 0.796 & 0.840\\
\bottomrule
\end{tabular}
\end{adjustbox}
\end{table}
The threshold-centred margin is not uniformly better. For atelectasis, whose
labelled cells are about 98\% positive and whose thresholds lie between 0.96
and 0.99, $|2p-1|$ ranks correctness well because confident positives are
almost always correct. The margin ranks it poorly.
\begin{table}[h]\centering\small
\caption{Comparator fits on the source calibration tier.}
\begin{adjustbox}{max width=\linewidth}\begin{tabular}{lrrrl}
\toprule
Model & Conformal $\alpha$ & Calibration acceptance & Expected-loss cutoff & Temperatures (Card., Edema, Cons., Atel., Eff.)\\
\midrule
M1 & 0.160 & 0.470 & -0.390 & 1.10, 1.13, 1.23, 0.95, 1.30\\
M2 & 0.270 & 0.591 & -0.471 & 1.28, 1.52, 1.77, 1.07, 1.58\\
M3 & 0.160 & 0.527 & -0.440 & 0.62, 0.66, 0.90, 0.98, 0.84\\
M4 & 0.150 & 0.566 & -0.311 & 1.24, 1.26, 1.41, 1.01, 1.47\\
\bottomrule
\end{tabular}
\end{adjustbox}
\end{table}

\section{Context interventions: selection versus prediction}
\small
\begin{longtable}{lllrr}
\toprule
Site/model/cond. & Term & & Estimate & 95\\% CI\\
\midrule
\endhead
external M2 C1 & policy effect & & +0.0333 & (+0.0305, +0.0360)\\
 & order1 prediction on C0 set & & +0.0307 & (+0.0278, +0.0333)\\
 & order1 selection & & +0.0027 & (+0.0016, +0.0037)\\
 & order2 selection & & +0.0202 & (+0.0190, +0.0214)\\
 & order2 prediction on intervention set & & +0.0131 & (+0.0102, +0.0158)\\
 & full cohort prediction effect & & +0.0131 & (+0.0102, +0.0158)\\
external M2 C2 & policy effect & & +0.0410 & (+0.0376, +0.0442)\\
 & order1 prediction on C0 set & & +0.0579 & (+0.0547, +0.0611)\\
 & order1 selection & & $-$0.0170 & ($-$0.0186, $-$0.0153)\\
 & order2 selection & & +0.0200 & (+0.0184, +0.0216)\\
 & order2 prediction on intervention set & & +0.0210 & (+0.0177, +0.0242)\\
 & full cohort prediction effect & & +0.0408 & (+0.0378, +0.0439)\\
external M3 C1 & policy effect & & $-$0.0134 & ($-$0.0157, $-$0.0112)\\
 & order1 prediction on C0 set & & $-$0.0127 & ($-$0.0145, $-$0.0108)\\
 & order1 selection & & $-$0.0007 & ($-$0.0023, +0.0009)\\
 & order2 selection & & +0.0063 & (+0.0048, +0.0080)\\
 & order2 prediction on intervention set & & $-$0.0197 & ($-$0.0215, $-$0.0179)\\
 & full cohort prediction effect & & $-$0.0159 & ($-$0.0176, $-$0.0141)\\
external M3 C2 & policy effect & & +0.0129 & (+0.0102, +0.0154)\\
 & order1 prediction on C0 set & & +0.0157 & (+0.0135, +0.0180)\\
 & order1 selection & & $-$0.0029 & ($-$0.0045, $-$0.0012)\\
 & order2 selection & & +0.0077 & (+0.0060, +0.0094)\\
 & order2 prediction on intervention set & & +0.0052 & (+0.0028, +0.0075)\\
 & full cohort prediction effect & & +0.0103 & (+0.0083, +0.0123)\\
external M4 C1 & policy effect & & +0.0092 & (+0.0069, +0.0116)\\
 & order1 prediction on C0 set & & +0.0132 & (+0.0109, +0.0153)\\
 & order1 selection & & $-$0.0040 & ($-$0.0054, $-$0.0025)\\
 & order2 selection & & +0.0011 & ($-$0.0002, +0.0025)\\
 & order2 prediction on intervention set & & +0.0081 & (+0.0059, +0.0103)\\
 & full cohort prediction effect & & +0.0108 & (+0.0088, +0.0128)\\
external M4 C2 & policy effect & & +0.0015 & ($-$0.0005, +0.0035)\\
 & order1 prediction on C0 set & & +0.0031 & (+0.0013, +0.0051)\\
 & order1 selection & & $-$0.0017 & ($-$0.0030, $-$0.0003)\\
 & order2 selection & & +0.0002 & ($-$0.0012, +0.0015)\\
 & order2 prediction on intervention set & & +0.0013 & ($-$0.0005, +0.0031)\\
 & full cohort prediction effect & & +0.0031 & (+0.0013, +0.0049)\\
source M2 C1 & policy effect & & +0.1523 & (+0.1384, +0.1669)\\
 & order1 prediction on C0 set & & +0.1479 & (+0.1336, +0.1626)\\
 & order1 selection & & +0.0044 & ($-$0.0002, +0.0092)\\
 & order2 selection & & +0.0281 & (+0.0232, +0.0329)\\
 & order2 prediction on intervention set & & +0.1242 & (+0.1093, +0.1397)\\
 & full cohort prediction effect & & +0.1242 & (+0.1093, +0.1397)\\
source M2 C2 & policy effect & & +0.1576 & (+0.1424, +0.1723)\\
 & order1 prediction on C0 set & & +0.1744 & (+0.1600, +0.1885)\\
 & order1 selection & & $-$0.0168 & ($-$0.0240, $-$0.0097)\\
 & order2 selection & & +0.0319 & (+0.0259, +0.0378)\\
 & order2 prediction on intervention set & & +0.1257 & (+0.1104, +0.1403)\\
 & full cohort prediction effect & & +0.1485 & (+0.1343, +0.1625)\\
source M3 C1 & policy effect & & $-$0.0182 & ($-$0.0259, $-$0.0096)\\
 & order1 prediction on C0 set & & $-$0.0181 & ($-$0.0250, $-$0.0110)\\
 & order1 selection & & $-$0.0001 & ($-$0.0051, +0.0047)\\
 & order2 selection & & +0.0075 & (+0.0025, +0.0128)\\
 & order2 prediction on intervention set & & $-$0.0257 & ($-$0.0330, $-$0.0178)\\
 & full cohort prediction effect & & $-$0.0223 & ($-$0.0288, $-$0.0152)\\
source M3 C2 & policy effect & & +0.0601 & (+0.0505, +0.0698)\\
 & order1 prediction on C0 set & & +0.0714 & (+0.0622, +0.0807)\\
 & order1 selection & & $-$0.0113 & ($-$0.0173, $-$0.0050)\\
 & order2 selection & & +0.0148 & (+0.0095, +0.0199)\\
 & order2 prediction on intervention set & & +0.0453 & (+0.0359, +0.0544)\\
 & full cohort prediction effect & & +0.0639 & (+0.0549, +0.0729)\\
source M4 C1 & policy effect & & +0.0003 & ($-$0.0077, +0.0090)\\
 & order1 prediction on C0 set & & +0.0000 & ($-$0.0070, +0.0074)\\
 & order1 selection & & +0.0003 & ($-$0.0050, +0.0056)\\
 & order2 selection & & +0.0054 & (+0.0003, +0.0106)\\
 & order2 prediction on intervention set & & $-$0.0051 & ($-$0.0121, +0.0024)\\
 & full cohort prediction effect & & $-$0.0004 & ($-$0.0071, +0.0065)\\
source M4 C2 & policy effect & & +0.0304 & (+0.0221, +0.0387)\\
 & order1 prediction on C0 set & & +0.0333 & (+0.0255, +0.0413)\\
 & order1 selection & & $-$0.0029 & ($-$0.0084, +0.0025)\\
 & order2 selection & & +0.0079 & (+0.0032, +0.0130)\\
 & order2 prediction on intervention set & & +0.0225 & (+0.0148, +0.0301)\\
 & full cohort prediction effect & & +0.0306 & (+0.0229, +0.0382)\\
\bottomrule
\end{longtable}

\normalsize

\section{C2 assignment audit}
\begin{table}[h]\centering\small
\caption{Replay of the C2 assignment (verified to reproduce \texttt{apply\_c2}
exactly).}
\begin{tabular}{lrr}
\toprule
Quantity & Source & External\\
\midrule
image rows & 16,456 & 135,561\\
studies & 14,766 & 132,923\\
rows with identical text after C2 & 6 & 34\\
\quad of which donated by a different patient & 6 & 34\\
rows with a same-patient donor & 0 & 0\\
rows donating to themselves & 0 & 0\\
multi-image studies & 1,607 & 2,608\\
\quad receiving context from $>1$ donor study & 1,465 & 2,564\\
\quad receiving $>1$ distinct context string & 1,459 & 2,403\\
rows whose target-term profile changed & 0.537 & 0.280\\
mean effective tokens & 13.047 & 7.304\\
\bottomrule
\end{tabular}

\end{table}

\section{Exploratory strata}
\begin{table}[h]\centering\small
\caption{View strata (corrected estimator; studies whose frontal images are
all AP or all PA).}
\begin{adjustbox}{max width=\linewidth}\begin{tabular}{llrrrrr}
\toprule
Stratum & Model & Studies src/ext & Risk src & Risk ext & Transfer gap (95\% CI)\\
\midrule
AP-only & M1 & 8,837/110,718 & 0.2056 & 0.2356 & +0.0300 (+0.0158, +0.0445)\\
 & M2 & 8,837/110,718 & 0.3076 & 0.2566 & $-$0.0510 ($-$0.0651, $-$0.0371)\\
 & M3 & 8,837/110,718 & 0.1872 & 0.1907 & +0.0035 ($-$0.0084, +0.0159)\\
 & M4 & 8,837/110,718 & 0.1366 & 0.1291 & $-$0.0075 ($-$0.0178, +0.0020)\\
\addlinespace
PA-only & M1 & 5,910/22,111 & 0.2264 & 0.1665 & $-$0.0599 ($-$0.0801, $-$0.0399)\\
 & M2 & 5,910/22,111 & 0.2947 & 0.4393 & +0.1446 (+0.1235, +0.1638)\\
 & M3 & 5,910/22,111 & 0.2129 & 0.1831 & $-$0.0298 ($-$0.0503, $-$0.0108)\\
 & M4 & 5,910/22,111 & 0.1465 & 0.1799 & +0.0335 (+0.0169, +0.0487)\\
\bottomrule
\end{tabular}
\end{adjustbox}
\end{table}
\begin{table}[h]\centering\small
\caption{External-site demographic strata (registered metrics MET016 and
MET017; no strata were registered; groups with fewer than 500 studies omitted;
max-of-groups statistics are biased upward). Source demographics were not
acquired.}
\begin{adjustbox}{max width=\linewidth}\begin{tabular}{llrrl}
\toprule
Variable & Model & Coverage disparity (MET016) & Worst-group risk (MET017) & Worst group\\
\midrule
sex & M3 & 0.004 (+0.000, +0.010) & 0.192 (+0.189, +0.197) & Female\\
 & M4 & 0.009 (+0.004, +0.014) & 0.138 (+0.136, +0.142) & Female\\
age band & M3 & 0.115 (+0.105, +0.125) & 0.195 (+0.192, +0.203) & 40-64\\
 & M4 & 0.033 (+0.024, +0.042) & 0.152 (+0.145, +0.159) & <40\\
race & M3 & 0.039 (+0.022, +0.064) & 0.214 (+0.205, +0.224) & Unknown\\
 & M4 & 0.007 (+0.006, +0.029) & 0.147 (+0.140, +0.168) & Pacific Islander\\
insurance type & M3 & 0.061 (+0.052, +0.070) & 0.196 (+0.192, +0.215) & Private Insurance\\
 & M4 & 0.019 (+0.013, +0.033) & 0.145 (+0.142, +0.156) & Medicaid\\
\bottomrule
\end{tabular}
\end{adjustbox}
\end{table}

\section{Training records}
The committed Stage 6b audit artifact reports an M2 curve whose epochs 1--4
(validation AUROC 0.50) belong to a discarded run with empty context input and
whose epochs 5--6 belong to the valid run. Its log parser kept the first
occurrence of each epoch number across logs. \Cref{tab:stage6} separates the
runs. The selected checkpoints are unaffected.
\begin{table}[h]\centering\small
\caption{Training runs recovered from the logs.}\label{tab:stage6}
\begin{adjustbox}{max width=\linewidth}\begin{tabular}{llrrl}
\toprule
Model & Log / run & Epochs & Best epoch (AUROC) & Validation macro AUROC by epoch\\
\midrule
M1 & stage6\_training.log / 0 & 6 & 3 (0.8774) & 0.868, 0.877, 0.877, 0.864, 0.864, 0.850\\
M2 & stage6\_training.log / 0 & 4 & 1 (0.5001) & 0.500, 0.500, 0.500, 0.500\\
M2 & stage6\_training\_m2m4.log / 0 & 6 & 3 (0.7629) & 0.757, 0.758, 0.763, 0.747, 0.729, 0.730\\
M4 & stage6\_training\_m2m4.log / 0 & 7 & 4 (0.8885) & 0.884, 0.888, 0.886, 0.888, 0.870, 0.874, 0.852\\
M1 & stage6\_seed2.log / 0 & 6 & 3 (0.8741) & 0.874, 0.870, 0.874, 0.863, 0.847, 0.853\\
M4 & stage6\_seed2.log / 0 & 3 & 1 (0.8850) & 0.885, 0.865, 0.884\\
M4 & stage6\_seed2.log / 1 & 2 & 1 (0.8876) & 0.888, 0.874\\
M4 & stage6\_seed2.log / 2 & 4 & 1 (0.8874) & 0.887, 0.873, 0.879, 0.877\\
M2 & revision\_gpu\_queue.log / 0 & 5 & 2 (0.7683) & 0.757, 0.768, 0.765, 0.756, 0.747\\
\bottomrule
\end{tabular}
\end{adjustbox}
\end{table}

\section{The 50\% boundary property}
Let $\hat\Delta\sim N(\Delta,s^2)$ and confirm when $\hat\Delta-1.96s>0$ and
$\hat\Delta\ge\delta$. Then
$\Pr(\text{confirm}\mid\Delta=\delta)=\Pr\{Z\ge\max(0,\,1.96-\delta/s)\}$. This
is at most $1/2$, equals $1/2$ when $\delta\ge1.96s$, and is smaller otherwise. The
property follows from the rule's construction and does not depend on the
application; we do not present it as a methodological contribution.

\section{Software and commands}
Python 3.14.4, NumPy 2.5.1, pandas 3.0.3, SciPy 1.18.0, scikit-learn 1.9.0,
PyTorch 2.13.0+cu126, torchvision 0.28.0, transformers 5.14.1, Pillow 12.3.0, on
a single 6\,GB GPU (GTX 1660 SUPER). Commands:
\begin{verbatim}
python protocols/C3E6_stage2/scripts/validate_protocol.py
python -m c3e.revision.runner        # all CPU analyses -> results/c3e_revision/
python -m c3e.revision.report        # tables, macros, figures -> paper/revision/
bash scripts/run_revision_gpu_queue.sh   # GPU steps (see below)
cd c3e_audit_toolkit && python -m pytest -q
\end{verbatim}

\section{GPU-dependent analyses}
\paragraph{Additional C2 donor permutations (post hoc).}
\begin{center}\small\begin{tabular}{llrrl}
\toprule
Site & C2 seed & Model & H4 (C2$-$C1), 95\% CI & Verdict\\
\midrule
external & 20260718 (frozen) & M2 & +0.0076 (+0.0049, +0.0102) & C\\
external & 20260718 (frozen) & M3 & +0.0263 (+0.0240, +0.0286) & C\\
external & 20260718 (frozen) & M4 & $-$0.0077 ($-$0.0101, $-$0.0054) & O\\
external & 20260719 & M2 & +0.0068 (+0.0042, +0.0095) & C\\
external & 20260719 & M3 & +0.0252 (+0.0228, +0.0275) & C\\
external & 20260719 & M4 & $-$0.0063 ($-$0.0086, $-$0.0039) & O\\
source & 20260718 (frozen) & M2 & +0.0053 ($-$0.0095, +0.0201) & I\\
source & 20260718 (frozen) & M3 & +0.0783 (+0.0691, +0.0865) & C\\
source & 20260718 (frozen) & M4 & +0.0301 (+0.0215, +0.0383) & C\\
source & 20260719 & M2 & $-$0.0140 ($-$0.0282, +0.0002) & I\\
source & 20260719 & M3 & +0.0779 (+0.0685, +0.0872) & C\\
source & 20260719 & M4 & +0.0274 (+0.0188, +0.0355) & C\\
source & 20260720 & M2 & $-$0.0069 ($-$0.0207, +0.0062) & I\\
source & 20260720 & M3 & +0.0801 (+0.0706, +0.0889) & C\\
source & 20260720 & M4 & +0.0335 (+0.0248, +0.0417) & C\\
\bottomrule
\end{tabular}\end{center}
Verdict codes as in the main text (C: interval above zero). Corrected estimator; all other inputs are the primary ones.
\paragraph{M2 and M3 under training seed 20260719 (post hoc).}
\begin{center}\small\begin{tabular}{lllrl}
\toprule
H & Model & Estimator & Estimate (95\% CI) & V\\
\midrule
H1 & M1 & original & +0.0511 (+0.0450, +0.0568) & C\\
H1 & M1 & evaluable-study & $-$0.0082 ($-$0.0194, +0.0027) & I\\
H1 & M2 & original & +0.0669 (+0.0586, +0.0749) & C\\
H1 & M2 & evaluable-study & $-$0.0279 ($-$0.0411, $-$0.0147) & O*\\
H1 & M3 & original & +0.0145 (+0.0086, +0.0208) & S\\
H1 & M3 & evaluable-study & $-$0.0402 ($-$0.0503, $-$0.0298) & O*\\
H1 & M4 & original & +0.0246 (+0.0182, +0.0315) & C\\
H1 & M4 & evaluable-study & $-$0.0469 ($-$0.0586, $-$0.0350) & O*\\
H2-C1 & M3 & original & +0.0154 (+0.0108, +0.0198) & S\\
H2-C1 & M3 & evaluable-study & +0.0150 (+0.0072, +0.0228) & S\\
H3 & M3 & original & $-$0.0366 ($-$0.0414, $-$0.0318) & O\\
H3 & M3 & evaluable-study & $-$0.0320 ($-$0.0405, $-$0.0237) & O\\
H4 source & M3 & original & +0.0356 (+0.0311, +0.0404) & C\\
H4 source & M3 & evaluable-study & +0.0565 (+0.0484, +0.0647) & C\\
H4 external & M3 & original & +0.0023 (+0.0006, +0.0040) & C\\
H4 external & M3 & evaluable-study & +0.0011 ($-$0.0010, +0.0031) & I\\
\bottomrule
\end{tabular}\end{center}


\section{Expert adjudication protocol (summary)}
A blinded adjudication of a stratified sample of studies at both sites is the
analysis that could arbitrate between label-handling alternatives. The full
protocol, annotation form, disagreement procedure and sample-size rationale are
in \texttt{docs/revision/EXPERT\_ADJUDICATION\_PROTOCOL.md}. It has not been
performed. In brief: two board-certified radiologists independently label the
five findings from the image alone, blinded to site, model output and report
labels, with a third resolving disagreements. Studies are sampled within strata
of site $\times$ label availability $\times$ model decision, and estimates are
weighted by inverse sampling fractions.

\end{document}
